% META: {"id":"1SPE-EXPO-CO-021","chapitre":"1SPE-EXPONENTIELLE","type_objet":"corrige","exercice_id":"1SPE-EXPO-EX-021","status":"approved","origin":"generated","status_history":["generated","audited","accepted","approved"],"release_acceptance":"RELEASE_OWNER_BATCH_ACCEPTANCE","acceptance_closure_digest":"sha256:9b3ccf9a81c5520fbb7e03b7b2d3e7bf2057a02834b6d908bf3be5f49f3b553f","acceptance_receipt_digest":"sha256:4fad86f0282e0fa4e0419cca1ad6379ae7af5a280c04a4a90880f90a1c044242"}
% BEGIN-VERIFY
% from sympy import *
% assert exp(3) < exp(5), "e^3 < e^5"
% assert exp(-2) < exp(-1), "e^{-2} < e^{-1}"
% assert exp(-3) < exp(0), "e^{-3} < e^0"
% END-VERIFY
\begin{corrige}{1SPE-EXPO-EX-021}
La fonction exponentielle est \textbf{strictement croissante} sur $\mathbb{R}$. Donc si $a < b$, alors $\mathrm{e}^a < \mathrm{e}^b$.

\textbf{1.} On a $3 < 5$, donc par croissance de l'exponentielle : $\mathrm{e}^3 < \mathrm{e}^5$.

\textbf{2.} On a $-2 < -1$, donc par croissance de l'exponentielle : $\mathrm{e}^{-2} < \mathrm{e}^{-1}$.

\textbf{3.} On a $-3 < 0$, donc par croissance de l'exponentielle : $\mathrm{e}^{-3} < \mathrm{e}^0 = 1$.
\end{corrige}
